\begin{lemma}
\label{lemma-cone-tor-amplitude}
Let $(X, \mathcal{O}_X)$ be a ringed space.
Let $(K, L, M, f, g, h)$ be a distinguished
triangle in $D(\mathcal{O}_X)$. Let $a, b \in \mathbf{Z}$.
\begin{enumerate}
\item If $K$ has tor-amplitude in $[a + 1, b + 1]$ and
$L$ has tor-amplitude in $[a, b]$ then $M$ has
tor-amplitude in $[a, b]$.
\item If $K$ and $M$ have tor-amplitude in $[a, b]$, then
$L$ has tor-amplitude in $[a, b]$.
\item If $L$ has tor-amplitude in $[a + 1, b + 1]$
and $M$ has tor-amplitude in $[a, b]$, then
$K$ has tor-amplitude in $[a + 1, b + 1]$.
\end{enumerate}
\end{lemma}

\begin{proof}
Omitted. Hint: This just follows from the long exact cohomology sequence
associated to a distinguished triangle and the fact that
$- \otimes_{\mathcal{O}_X}^{\mathbf{L}} \mathcal{F}$
preserves distinguished triangles.
The easiest one to prove is (2) and the others follow from it by
translation.
\end{proof}
